\chapter{Hatha Texts, Diet and Healthy Living}
\label{ch:hatha-diet-living}

\emph{What do the hatha texts teach, what do they and the Gita say about food and daily living, and why does the practitioner's own experience matter?}
This chapter completes the texts unit. It covers the Hatha Pradipika and, beside it, the Gheranda Samhita; the two ideas of diet, mitahara and yuktahara; the Gita's picture of a healthy life; and the place of subjective experience in daily practice. The Yoga Sutra, the Gita's selected slokas and the Prasthanatrayi are in \cref{ch:texts}.

\begin{unverified}
This chapter is written from notes blocks whose fact check is still pending. Points follow the official guidebook's chapters on hatha yoga and health, the 1915 English translation of the Hatha Pradipika on Wikisource, and the Gita, as marked in the notes. Verse numbers are to be checked against a standard edition. A reviewer must check it before print.
\end{unverified}

\section{The Hatha Pradipika}

\subsection*{The text and its idea}
The Hatha Pradipika is a classical Sanskrit manual of hatha yoga by Svatmarama, a medieval text commonly placed in the fifteenth century; the guidebook calls Svatmarama a disciple of Gorakhnath. Sources differ on his dates, so ``medieval'' is the safe word.

Hatha yoga is a body-based path that uses posture and breath to purify the body and balance body and mind, so that the body can support deeper meditation. It is body-oriented but not only for the body: the guidebook stresses that hatha and raja yoga are one whole, the first preparing the way for the second. The syllables \emph{ha} and \emph{tha} are read as the sun and the moon, the balance of the two energies carried by the pingala and ida channels. A second reading of the word is ``forceful effort''. The Hatha Pradipika and the Gheranda Samhita are the two classical manuals.

\subsection*{Four steps and four chapters}
The text's yoga has four steps: asana, pranayama, mudra and bandha, and \emph{nādānusandhāna}, attention to the inner sound. Its four chapters follow them:
\begin{itemize}
\item \textbf{The first chapter}: asana, with measured diet (mitahara), the helps and hindrances on the path, and the ten yamas and ten niyamas of the text; the opening verses praise hatha yoga and describe a suitable place for practice.
\item \textbf{The second chapter}: pranayama, beginning with the cleansing of the nadis (channels), the six cleansing actions, and the eight kumbhakas with their cautions. Pranayama is shown as a way to calm the mind, balance prana and support meditation; the text warns that improper practice brings harm.
\item \textbf{The third chapter}: the mudras and bandhas and the rousing of kundalini. It names ten practices: maha mudra, maha bandha, maha vedha, khechari, uddiyana, mula and jalandhara bandha, viparita karani, vajroli and shakti chalana.
\item \textbf{The fourth chapter}: nadanusandhana and samadhi, with the inner sound and its stages, laya (absorption of the mind through sound), and four stages of progress, arambha (beginning), ghata, parichaya (familiarity) and nishpatti (completion).
\end{itemize}

\subsection*{Asanas}
The text says Shiva taught eighty-four asanas, and it describes fifteen. Four are named as the important ones: siddhasana, padmasana, simhasana and bhadrasana, and siddhasana is called the most comfortable. Others described include svastikasana, gomukhasana, virasana, kurmasana, kukkutasana, dhanurasana, matsyasana, paschimottanasana, mayurasana and shavasana.

\subsection*{Six helps and six hindrances}
The helps to success (\emph{sādhaka tattva}):
\begin{itemize}
\item \textbf{Utsaha}: enthusiasm, a positive attitude.
\item \textbf{Sahasa}: courage to meet inner experience.
\item \textbf{Dhairya}: patience and steadiness in practice.
\item \textbf{Tattvajnana}: knowledge that ties actions to the goal.
\item \textbf{Nishchaya}: firm resolve.
\item \textbf{Janasanga-parityaga}: keeping away from crowds and social distraction.
\end{itemize}
The hindrances (\emph{bādhaka tattva}):
\begin{itemize}
\item \textbf{Atyahara}: overeating.
\item \textbf{Prayasa}: overexertion.
\item \textbf{Prajalpa}: too much talk.
\item \textbf{Niyamagraha}: rigid clinging to rules, such as fixed habits of cold baths or strict diets.
\item \textbf{Janasanga}: the company of crowds.
\item \textbf{Laulya}: fickleness, an unsteady and restless mind.
\end{itemize}

\subsection*{Cleansing and pranayama in the text}
The six cleansing actions are dhauti, basti, neti, trataka, nauli and kapalabhati. The text advises them first for those with excess fat or phlegm, and says that others do not need them. Its pranayama is nadi shodhana and eight kumbhakas: surya bhedana, ujjayi, sitkari, sitali, bhastrika, bhramari, murcha and plavini. The guidebook's chapter summary speaks of ten types but lists eight; teach the eight with nadi shodhana.

\subsection*{Why the text matters}
It is the base of later hatha practice. It joins bodily and spiritual practice, gives step-by-step guidance, and stresses discipline, perseverance and correct practice.

\section{The Gheranda Samhita}

The Gheranda Samhita teaches a seven-limbed yoga, called \emph{saptāṅga} or \emph{ghaṭastha} yoga, in seven chapters, one for each limb. Its image is the body and mind as a clay vessel (\emph{ghaṭa}) that must be fired and made firm by yoga. \emph{Ghaṭa śuddhi}, the purification of the whole body-mind structure, is what the seven limbs bring about, and its aim is self-realisation.

\textit{The chapter contents below come from the workshop slides and are to be checked against a standard edition.}
\begin{itemize}
\item \textbf{Shatkarma}: the six cleansing actions, done before the other limbs.
\item \textbf{Asana}: thirty-two postures, grouped as meditative, relaxing and body-strengthening, with steadiness and comfort stressed.
\item \textbf{Mudra}: twenty-five mudras that direct prana, for example maha mudra, nabho mudra, khechari, viparita karani and yoni mudra.
\item \textbf{Pratyahara}: withdrawal of the senses, as a preparation for meditation.
\item \textbf{Pranayama}: advice on diet first, then eight kumbhakas (sahita, surya bheda, ujjayi, sitali, bhastrika, bhramari, murchha and kevali), with signs of success.
\item \textbf{Dhyana}: three kinds of meditation, sthula (gross, on a form), jyotis (on light) and sukshma (subtle).
\item \textbf{Samadhi}: six kinds, through dhyana yoga, nada yoga, rasananda, laya siddhi, bhakti yoga and raja yoga; the aim is liberation and knowledge of the Self.
\end{itemize}
Other hatha texts include the Hatharatnavali and the Shiva Samhita.

\section{Diet: mitahara and yuktahara}

The yogic idea of diet is about quality and the right amount, not calories: moderate, wholesome food and a moderate way of living. It rests on two terms, one from each text.

\subsection*{Mitahara in the Hatha Pradipika}
\emph{Mita} means limited or measured. Mitahara is food that is nourishing, pleasant, easy to digest and liked, eaten in the right amount: about three parts of the hunger, leaving the fourth part of the stomach empty. The food is soft or lightly oiled rather than dry, and naturally sweet. The text speaks of eating as an offering; read it as eating with attention and thankfulness.

The text lists suitable (\emph{pathya}) food: wheat, rice, barley, milk, ghee, butter, sugar candy, honey, dry ginger, green gram and certain leafy vegetables. Another verse calls a diet with too many vegetables unwholesome, so say fresh vegetables in moderation. It lists unsuitable (\emph{apathya}) food: food that is too bitter, sour, pungent, salty or hot; stale, reheated or dry; heavy and hard to digest; and, by name, oil, mustard, sesame, curds, buttermilk, garlic, alcohol, fish and meat. Overeating is itself unsuitable.

\subsection*{Yuktahara in the Bhagavad Gita}
The Gita (sixth chapter, verse 17) says that yoga ends sorrow for the one who is moderate in eating, recreation, effort, sleep and waking. Yuktahara is therefore wider than food: balance in food and in activity, rest and sleep. The guidebook adds that \emph{āhāra} covers what is taken in through the senses and the mind as well as through the mouth.

\subsection*{The three kinds of food}
The Gita's seventeenth chapter sorts food by the gunas:
\begin{itemize}
\item \textbf{Sattvic} food gives life, purity, strength, health, joy and cheer; it is savoury, bland, substantial and pleasing. It supports clarity and calm.
\item \textbf{Rajasic} food is stimulating: very spicy, bitter, sour, salty or hot.
\item \textbf{Tamasic} food is stale, spoiled, left over or impure, and weighs the mind down. The guidebook counts over-processed food as tamasic.
\end{itemize}

\subsection*{In practice}
A sattvic diet is fresh, simple food: grains, fruit, vegetables, pulses, nuts, milk and ghee. Change the diet gradually. Ayurveda names six tastes: sweet, sour, salty, pungent, bitter and astringent; the yogic diet leans to the sweet taste, in the sense of naturally pleasant and nourishing food. Vajrasana is the one posture commonly allowed soon after a meal; other practice is done on an empty stomach.

The lists of suitable and unsuitable food differ between texts, and with the person's climate and health. A vegetarian diet is the tradition; say so as the tradition, not as a rule to impose on students.

\section{Healthy living in the Bhagavad Gita}

The Gita is a source of ideas on daily conduct and on the mind. An exam answer gives the idea and its verse; it makes no medical claim.
\begin{itemize}
\item \textbf{Regulated living} (sixth chapter): moderation in food, recreation, effort, sleep and waking.
\item \textbf{The chain to anger} (second chapter): dwelling on the objects of the senses leads to attachment, attachment to desire, desire to anger, and anger to confusion and the loss of memory and reason. It is the Gita's account of the stress cycle.
\item \textbf{Even-minded action}: steady action without attachment, without liking and disliking, and without hunger for the fruit; the selected verses of the second chapter are in \cref{ch:texts}.
\item \textbf{The gunas in daily life} (fourteenth and seventeenth chapters): sattva binds to happiness and knowledge, rajas to action, tamas to heedlessness and sleep; food, work and the mind follow the same three.
\end{itemize}

\subsection*{The divine qualities}
The sixteenth chapter lists the qualities of a person of divine nature (\emph{daivī sampat}), the base of a healthy life in society. The first fourteen: fearlessness (abhaya), purity of mind (sattva-samshuddhi), steadiness in knowledge (jnana-yoga-vyavasthiti), charity (dana), control of the senses (dama), sacrifice (yajna), self-study (svadhyaya), austerity (tapas), straightforwardness (arjava), non-violence (ahimsa), truthfulness (satya), freedom from anger (akrodha), renunciation (tyaga) and peace (shanti). Twelve more follow: not finding fault (apaishuna), compassion (daya), freedom from greed (aloluptva), gentleness (mardava), modesty (hri), freedom from restlessness (achapala), vigour (tejas), forgiveness (kshama), fortitude (dhriti), cleanliness (shaucha), freedom from hatred (adroha) and freedom from pride (na-atimanita).

In short, healthy living in the Gita means balanced food, sleep, work and rest; a calm and steady mind; and right action without strain. The workshop adds living in step with natural rhythms and avoiding both indulgence and harsh self-denial; daily meditation and mindfulness are teaching summaries, not direct statements of the text. Contentment, santosha, is a niyama of the Yoga Sutra; the Gita values the same attitude in acting without attachment.

\section{Subjective experience in daily practice}

Subjective experience is what the practitioner feels and notices in the body, breath and mind, during and after practice. It is the practitioner's own report, not a measurement. Yoga is learned by doing, with steady attention, and the Yoga Sutra says that practice becomes firm when it is long, unbroken and done with care (\emph{abhyāsa}). Experience is how the practitioner knows whether the practice is working.

Before and after practice, note the breath, the body, the mood and the energy. Notice what is easier or harder than before. Adjust the practice to your own capacity, and tell the teacher what you notice. Look for ease and steadiness in the posture, a smooth breath, a calmer mind and less tension, and for the signs to stop: pain, dizziness or strain. One's feelings are a guide, not a diagnosis; a persistent problem needs a doctor.

Several ideas from the Yoga Sutra fit here. Svadhyaya, self-study, builds understanding of oneself and of the texts. Ishvara pranidhana, surrender, the guidebook explains as surrendering one's ego; it is commonly read as practising without clinging to the result. \emph{Eka-tattva abhyāsa}, practice on one principle, is the Sutra's remedy for a scattered mind. The workshop adds three modern words for the same habit: self-awareness, mindfulness and the mind-body connection. A short daily note of the breath and the mood is a simple way to start.

\begin{recall}
\begin{itemize}
\item Hatha Pradipika: Svatmarama; four steps, asana, pranayama, mudra and bandha, nadanusandhana; fifteen asanas described.
\item Six helps and six hindrances; six cleansing actions; eight kumbhakas.
\item Gheranda Samhita: seven limbs, shatkarma, asana, mudra, pratyahara, pranayama, dhyana, samadhi; ghata shuddhi.
\item Mitahara (Hatha Pradipika): measured food, a quarter of the stomach empty. Yuktahara (Gita): moderation in food, recreation, effort, sleep and waking.
\item Food: sattvic, rajasic, tamasic. Divine qualities: Gita, sixteenth chapter.
\end{itemize}
\end{recall}

\subsection*{Worked examples}
\begin{itemize}
\item \textbf{Model answer: compare mitahara and yuktahara.} Mitahara is the hatha term for measured, pure and wholesome food eaten with attention, leaving a quarter of the stomach empty. Yuktahara is the Gita's wider idea of balance in food and also in recreation, effort, sleep and waking.
\item \textbf{A student fasts for long hours before class to feel ``light''.} Explain that the aim is a light stomach, not an empty body: measured, regular meals, with practice on an empty or light stomach. Long fasting can bring dizziness and is not mitahara.
\item \textbf{A student pushes through pain to keep up with the class.} Use subjective experience: pain, dizziness and strain are signs to stop; adjust to your own capacity and tell the teacher.
\end{itemize}

\subsection*{Check yourself}
\begin{enumerate}
\item Who wrote the Hatha Pradipika?
\item Name its four steps.
\item Which four asanas does it name as the important ones?
\item Name the six hindrances in the text.
\item Name the seven limbs of the Gheranda Samhita.
\item What does ghata mean in the Gheranda Samhita?
\item What does mita mean in mitahara, and how much of the stomach is left empty?
\item Which text gives yuktahara, and in which chapter?
\item Name the three kinds of food in the Gita.
\item Which chapter of the Gita lists the divine qualities? Name three.
\item Why is subjective experience important in daily practice?
\item What does eka-tattva abhyasa mean?
\end{enumerate}

\subsection*{Answers}
\begin{enumerate}
\item Svatmarama.
\item Asana, pranayama, mudra and bandha, nadanusandhana.
\item Siddhasana, padmasana, simhasana and bhadrasana.
\item Overeating, overexertion, too much talk, rigid rules, the company of crowds, fickleness.
\item Shatkarma, asana, mudra, pratyahara, pranayama, dhyana, samadhi.
\item A vessel: the body and mind to be purified and made firm by yoga.
\item Limited or measured; about a quarter.
\item The Bhagavad Gita, sixth chapter.
\item Sattvic, rajasic and tamasic.
\item The sixteenth; for example fearlessness, truthfulness and compassion.
\item It lets the practitioner notice the effect of practice and adjust it to their own capacity.
\item Practice on one principle, to steady a scattered mind.
\end{enumerate}

\paragraph{Sources.} Hatha Pradipika of Svatmarama (English translation on Wikisource, 1915); the official guidebook for Level I (2016, older scheme), chapters on hatha yoga, Patanjali and health; the Bhagavad Gita (sixth, fourteenth, sixteenth and seventeenth chapters); the Gheranda Samhita (to be checked in a standard edition); the Board syllabus for Level 1.
